\conclusion
В работе рассмотрены теоретические и прикладные аспекты МКВ и составлен план прототипа защищенного МО.

Теоретические основы (раздел~\ref{sec:theoretical-foundations}) задают модели нарушителя --- получестную, злонамеренную и скрытую, --- допущение о честном большинстве и свойства корректности, приватности и справедливости. Криптографические примитивы (раздел~\ref{sec:cryptographic-foundations}) и классификация протоколов (раздел~\ref{sec:protocol-classification}) сопоставлены в таблице~\ref{tab:mpc-primitives}: выбор примитива определяется числом участников, моделью нарушителя и допустимыми издержками.

Анализ применений (раздел~\ref{sec:applications}) показал, что наиболее разработанным направлением является МО на разделенных данных. Существующие решения (SecureML, ABY3, TFEncrypted, CrypTen, Cerebro) различаются числом серверов и допущениями о нарушителе и в основном рассчитаны на получестную модель. Основными ограничениями остаются коммуникационные издержки, аппроксимация нелинейных функций, усечение чисел с фиксированной запятой, утечки через результат вычисления и доверие к среде исполнения (раздел~\ref{sec:limitations}); их устранение требует отдельных решений, таких как дифференциальная приватность и доверенные среды выполнения.

Для будущего исследования предложен продукт из двух раздельных блоков: модели логистической регрессии и заменяемого МКВ-модуля (раздел~\ref{sec:future-research}). Модуль строится на TFEncrypted и ABY3 на трех серверах в получестной модели при честном большинстве. Цель MVP --- показать возможность безопасной реализации совместного обучения и инференса: входы, веса и градиенты не восстанавливаются ни одним участником и координатором, а результат отличается от открытой контрольной реализации в пределах заданных допусков (см. подраздел~\ref{subsec:mvp-implementation-plan}). Реализация и измерения предстоят. Демонстрация не заменяет доказательство безопасности; защита от злонамеренного нарушителя, гарантии справедливости и промышленная готовность не заявляются.

После MVP предусматриваются расширение испытаний и отдельное исследование дифференциальной приватности, доверенных сред выполнения, федеративного обучения и ORAM (см. подраздел~\ref{subsec:mvp-risks-roadmap}). Замена МКВ-модуля при сохранении границы блоков потребует повторной проверки численной совместимости и модели угроз.